\section{Legal Properties and Limitations}
\label{sec:legal}

Simulated Annealing Bisection is a \emph{structure-layer} algorithm: it
changes how the bisection tree nodes are solved, not what the algorithm
optimises.
Its legal character is therefore inherited from the bisection tree structure
of \ar{} \citep{dellaLibera2026apportion}.
SA refines \metis{} solutions using only census-tract adjacency and
population data; no partisan data enters the SA loop.

\paragraph{Compactness under state law.}
SA's primary contribution is compactness improvement.
Polsby-Popper scores are recognised as one measure of compactness in
redistricting statutes and case law.
However, PP improvement alone is not a constitutional standard: courts apply
multiple compactness measures and weigh compactness against other
redistricting criteria (contiguity, equal population, community of interest
preservation).
A plan produced by SA should report PP alongside Reock and any
state-mandated compactness measures.

\paragraph{Post-\textit{Rucho} context.}
Under \textit{Rucho v.\ Common Cause} (2019), federal courts have no
jurisdiction over partisan gerrymandering claims; state courts apply state
constitutional standards.
SA does not alter the partisan character of the bisection tree structure and
produces consistently proportional outcomes in our three-state evaluation
(Table~\ref{tab:main}), consistent with B.17's finding that structure
dominates search.
This proportionality is a consequence of the prime-factor bisection
geometry, not of SA itself.

\paragraph{VRA considerations.}
SA's boundary-level adjustments operate within a fixed district count $k$
and bisection topology.
SA does not guarantee preservation of majority-minority district
configurations; VRA compliance must be verified separately using
\texttt{redist analyze --types bloc-voting} after plan generation
\citep{dellaLibera2026callais}.

\paragraph{Audit chain.}
Each SA run records the per-node seeds (Definition~\ref{def:sa-seed}),
cooling schedule parameters ($T_0$, $T_{\mathrm{final}}$,
\texttt{steps\_per\_tract}), the initial \metis{} \ec{}, the final
best-ever \ec{}, and the total steps run.
These fields satisfy the reproducibility requirements of the \dia{} seed
formula (prefix \texttt{"DIA\_SEED\_V1"}) and enable a verifier to
reproduce the temperature schedule and confirm that the submitted plan was
reachable under the stated parameters.
